\chapter*{Resumen}
\addcontentsline{toc}{chapter}{Resumen}
\label{pre:resumen}

\begingroup
\setstretch{1.1}
\setlength{\parindent}{0pt}
\setlength{\parskip}{0.7em}

\textbf{Título:} \TituloTrabajo\footnote{Trabajo de grado, modalidad Trabajo de Investigación.}

\textbf{Autores:} \AutorUno; \AutorDos.\footnote{\Facultad. \Escuela. Director: \DirectorNombre, \DirectorTitulo Codirector: \CodirectorNombre, \CodirectorTitulo}

\textbf{Palabras clave:} muografía por transmisión, exploración volcánica, aprendizaje profundo, simulación Monte Carlo, mejora de imágenes, baja estadística.

\textbf{Descripción:}

% Cifras tomadas de los capítulos 4 y 5 (evaluación final EF-01 con la corrección C2).
La muografía por transmisión permite estimar la distribución de densidad en el interior de un volcán a partir de la atenuación del flujo de muones cósmicos, pero el bajo flujo natural obliga a exposiciones de semanas o meses para que los contrastes de densidad se distingan del ruido de conteo. Este trabajo desarrolló y evaluó un esquema de mejora de imágenes muográficas de baja estadística basado en aprendizaje profundo y entrenado con datos sintéticos. Mediante simulación Monte Carlo con Geant4, con una fuente de muones derivada de una parametrización del flujo y del generador EcoMug, se obtuvieron 114 imágenes de referencia de diecinueve escenarios volcánicos idealizados y, por adelgazamiento binomial, 9\,120 pares de entrada y referencia con el 50, el 10, el 5 y el 1\,\% de los conteos, particionados por geometría. Con el 5\,\% de los conteos se diseñó, entrenó y evaluó una red generativa antagónica condicional de Wasserstein con penalización de gradiente, formada por un generador U-Net y un crítico PatchGAN, junto con una U-Net entrenada con pérdida L1 y variantes ajustadas para conservar las cuentas. En la evaluación final, descriptiva y sobre datos potencialmente expuestos, porque las escenas reservadas ya se habían evaluado antes, la red generativa aumentó en prueba el SSIM de 0,7643 a 0,8342 y el PSNR de 24,258 a 31,170~dB respecto de su entrada y, según ambas métricas, superó en cada una de las cuatro geometrías reservadas a la imagen sin procesar con el doble de conteos, aunque no a la obtenida con diez veces más. Al mismo tiempo, aumentó el sesgo regional de las cuentas de 1,832 a 14,566\,\%, y las variantes que lo redujeron disminuyeron la similitud. La calidad comparable a la de alta estadística quedó, por ello, respaldada solo de forma parcial, y el factor de reducción del tiempo de adquisición alcanzable no quedó demostrado: el valor nominal de 20 describe solo la reducción de conteos de la entrada. Los resultados se limitan a escenarios sintéticos idealizados y a una sola fracción de conteo; su extensión a datos de campo exige un criterio de calidad con un margen fundamentado y la respuesta de un detector físico.

\endgroup
